% BHU PYQ -- Professional Exam Template
\documentclass[11pt,a4paper]{article}

\usepackage[a4paper,top=2.5cm,bottom=2.5cm,left=2.8cm,right=2.8cm,headheight=14pt]{geometry}
\usepackage{amsmath,amssymb}
\usepackage[T1]{fontenc}
\usepackage[utf8]{inputenc}
\usepackage{lmodern}
\usepackage{microtype}
\usepackage{fancyhdr}
\usepackage{enumitem}
\usepackage{listings}
\usepackage{hyperref}
\usepackage{lastpage}
\usepackage{parskip}

\hypersetup{colorlinks=true,urlcolor=black,linkcolor=black,
  pdftitle={BVCA-104: Skill Training in R -- BHU PYQ}}

\pagestyle{fancy}
\fancyhf{}
\renewcommand{\headrulewidth}{0.4pt}
\renewcommand{\footrulewidth}{0.4pt}
\fancyhead[L]{\small\textit{Banaras Hindu University}}
\fancyhead[R]{\small\textit{BVCA-104: Skill Training in R}}
\fancyfoot[C]{\small Page \thepage\ of \pageref{LastPage}}

\lstset{
  basicstyle=\small\ttfamily,
  frame=single,
  breaklines=true,
  columns=flexible,
  keepspaces=true
}

\newlist{parts}{enumerate}{2}
\setlist[parts,1]{label=(\alph*),leftmargin=*,itemsep=4pt,topsep=3pt}
\setlist[parts,2]{label=(\roman*),leftmargin=*,itemsep=2pt,topsep=2pt}
\newcommand{\pts}[1]{\hfill\textit{[#1]}}

\begin{document}

\begin{center}
  {\large\bfseries BANARAS HINDU UNIVERSITY}\\[2pt]
  {\normalsize B.Voc. Semester I Examination, 2022-23}\\[6pt]
  \rule{0.8\linewidth}{0.5pt}\\[6pt]
  {\large\bfseries Paper: BVCA-104 --- Skill Training in R}\\[4pt]
  {\normalsize Time: Three Hours \hfill Maximum Marks: 70}
\end{center}

\rule{\linewidth}{0.4pt}

\smallskip
\noindent\textit{Instructions: Answer all Sections. Candidates are required to write their answers in their own words as far as practicable.}

\bigskip

\bigskip\noindent{\large\bfseries SECTION --- A: Long Answer Questions}\pts{$15 \times 2 = 30$}
\rule{\linewidth}{0.4pt}\smallskip

\noindent\textit{Instructions: Long answer type questions to be answered in about 750 words each.}

\medskip\noindent\textbf{Question 1.} List out the different control structures in R Programming with general syntaxes and brief explanation. Write R code for the function using if-else command:
\[
f(x, y) = \begin{cases}
x + y & \text{if } x < y \\
1 - x - y & \text{if } x > y \\
0 & \text{if } x = y
\end{cases}
\]
\centerline{\textbf{OR}}

\medskip\noindent\textbf{Question 2.} Define nested if-else and nested loop statements with general syntaxes and explanation. Write R codes for finding a maximum number from the given three numbers.

\medskip\noindent\textbf{Question 3.} Write a brief note on R. Write R codes for computing matrix inverse for a $3 \times 3$ matrix.
\centerline{\textbf{OR}}

\medskip\noindent\textbf{Question 4.} Illustrate arithmetic, logical, and relational operators with an example for each. Write R program for the factorial function using a for-loop.

\bigskip\noindent{\large\bfseries SECTION --- B: Short Answer Questions}\pts{$5 \times 6 = 30$}
\rule{\linewidth}{0.4pt}\smallskip

\noindent\textit{Instructions: Short answer type questions to be answered in about 550 words each. Attempt any six.}

\medskip\noindent\textbf{Question 1.} Write general syntax of R functions that are used to create bar and pie charts. Show the percentage of students registered in Data Science, R Programming, and Machine Learning subjects using these charts. The percentages are 30\%, 40\%, and 30\% in Data Science, R Programming, and Machine Learning, respectively.

\medskip\noindent\textbf{Question 2.} Giving a suitable example for each, explain the uses of the following operators:
\[
\text{"\%x\%"}, \quad \text{"\%*\%"}, \quad \text{"\%\%"}, \quad \text{"\%/\%"}, \quad \text{and} \quad \text{"\%in\%"}
\]

\medskip\noindent\textbf{Question 3.} Using \texttt{function()} command, write R codes to enter the following mathematical expressions in R:
\begin{parts}
  \item $y = 10^x + \sqrt{x}\log_{10}(x)$
  \item $y = e^{x/2} - \sin(x)\tan(|x|^4)$
\end{parts}

\medskip\noindent\textbf{Question 4.} State important features of R.

\medskip\noindent\textbf{Question 5.} Differentiate between vector, list, matrix, and data frame with general syntaxes and examples.

\medskip\noindent\textbf{Question 6.} Write all R codes to perform correlation and regression analysis for the following data set:
\begin{center}
\begin{tabular}{|c|c|c|c|c|c|c|}
\hline
\textbf{Height (inches)} & 60 & 62 & 65 & 68 & 70 & 72 \\ \hline
\textbf{Weight (lbs)} & 110 & 120 & 140 & 160 & 170 & 190 \\ \hline
\end{tabular}
\end{center}

\medskip\noindent\textbf{Question 7.} Write R syntaxes to apply t-test (single sample and two sample) and Chi-square test. For illustration, you may take data of your choice.

\medskip\noindent\textbf{Question 8.} Explain the use of \texttt{next}, \texttt{break}, and \texttt{stop} functions. Give an example for each.

\bigskip\noindent{\large\bfseries SECTION --- C: Very Short Answer Questions}\pts{$10 \times 1 = 10$}
\rule{\linewidth}{0.4pt}\smallskip

\noindent\textit{Instructions: Answer the following questions in a single word or sentence.}

\begin{enumerate}[label=\arabic*.,leftmargin=*,itemsep=6pt]
  \item How do you enter integer, complex, character, and numeric data in R?
  \item Write R code to enter the following data in R using \texttt{seq()} and \texttt{rep()}:
  \[
  1.0, \, 1.0, \, 1.0, \, 1.0, \, 1.0, \, 0.5, \, 1.0, \, 1.5, \, 2.0, \, 2.5, \, 3.0
  \]
  \item Write general syntax of the import/export of a CSV file.
  \item Write a code to compute sum, mean, and product of a given vector of elements.
  \item Write R code to add a legend in R plot.
  \item How do you set the margin of the plot in R?
  \item R is an implementation of \underline{\hspace{1.5cm}} language.
  \item Write R code to extract a column of a data frame.
  \item Write down the output of the codes:
  \begin{lstlisting}[language=R]
N = 1:15
which(N >= 6 & N < 10)
  \end{lstlisting}
  \item Write down the output of the codes:
  \begin{lstlisting}[language=R]
x = letters[1:10]
x[seq(2, 10, by = 2)]
  \end{lstlisting}
\end{enumerate}

\vfill
\rule{\linewidth}{0.4pt}
\begin{center}\small
\href{https://bhu-exam.vercel.app/}{Click here to visit our website}\\[4pt]\href{https://me-aryan.vercel.app/}{Click here to visit our contributor}\\[4pt]\href{https://quashan-me.vercel.app/}{Click here to visit our contributor}
\end{center}

\end{document}
